% !PRINTED 602
\weberpage{602}

Die Konstante $g$ ist nach (5) näher bestimmt:

\weq{7}{g = \frac{\gamma\,\pi}{M^2\sqrt{-\Delta}},}

wenn $\gamma$ die Anzahl der in $A$ enthaltenen Zahlklassen nach dem Modul $M$ ist.

Will man unter $T$ die Anzahl der ganzen Ideale der Hauptklasse $\overline{E}A$ verstehen, deren Norm nicht größer als $t$ ist, so bleibt (6) auch dann noch bestehen, nur hat man, wenn $-1$ in $A$ vorkommt, die Zahl $g$ zu halbieren, und in den beiden Ausnahmefällen $\Delta = -4$, $\Delta = -3$ möglicherweise noch durch 4 oder durch 6 zu teilen.

\begin{center}
\textbf{\S\,163.\qquad \gesperrt{Die Dirichletschen Summen.}}
\end{center}

Wir verstehen jetzt unter $\overline{O}$ die Idealgruppe der sämtlichen Ideale in $\Omega$ (mit Rücksicht auf den Exkludenten $S$); $A$ sei eine Kongruenzgruppe und

\weq{1}{\overline{A} = \overline{E}A,}

die Gruppe der entsprechenden Hauptideale. Wir zerlegen $\overline{O}$ nach \S\,161, (3) in die \gesperrt{Klassen} nach $A$:

\weq{2}{\overline{O} = \overline{A}_1 + \overline{A}_2 + \cdots + \overline{A}_h,}

und betrachten die Summen:

\weq{3}{S_i = \sum^{\mathfrak{a}_i}\frac{1}{(N\mathfrak{a}_i)^s},}

worin $s$ eine positive Variable $> 1$ bedeutet, und $\mathfrak{a}_i$ die sämtlichen ganzen Ideale einer Klasse $\overline{A}_i$ durchläuft.

Die Summen $S_i$ formen wir in folgender Weise um:

Da $N(\mathfrak{a}_i)$ eine der Zahlen 1, 2, 3, \dots\ sein muß, so bezeichnen wir mit $a_\nu$ die Anzahl der Ideale der Klasse $\overline{A}_i$, deren Norm $= \nu$ ist, und erhalten

\weq{4}{S_i = a_1 + \frac{a_2}{2^s} + \frac{a_3}{3^s} + \frac{a_4}{4^s} + \cdots}

wobei, wenn eine Zahl $\nu$ unter der $N(\mathfrak{a}_i)$ nicht vorkommt, $a_\nu = 0$ zu setzen ist.

Bezeichnen wir mit $Z(\nu)$ die Anzahl der Ideale der Klasse $\overline{A}_i$, deren Norm nicht größer als $\nu$ ist, so erhalten wir für $Z(\nu)$ aus der Formel (6) des vorigen Paragraphen einen Grenzwert:
%CONT>>
